\documentclass[10pt,a4paper]{amsart}
\usepackage[margin=19mm]{geometry}
\usepackage{amsmath,amssymb,mathtools}
\usepackage[scr=rsfs,cal=euler]{mathalfa}
\usepackage{xfrac}
\usepackage[unicode,hidelinks,bookmarks=false]{hyperref}
\newcommand{\ie}{i.e.\ }
\newcommand{\cf}{cf.\ }
\newcommand{\resp}{respectively\ }
\newcommand{\Subsection}{\subsection}
\providecommand{\nobreakdash}{\nobreak\hbox{-}}
\setlength{\emergencystretch}{2em}
\multlinegap0pt
\usepackage{bm}
\usepackage{bbm}
\usepackage{dsfont}
\usepackage{xspace}
\usepackage{cancel}
\usepackage{mathtools}
\usepackage[shortlabels]{enumitem}
\usepackage{amsthm}
\makeatletter
\@ifundefined{thm}{\newtheorem{thm}{Thm}}{}
\@ifundefined{claim}{\newtheorem{claim}[thm]{Claim}}{}
\@ifundefined{fact}{\newtheorem{fact}[thm]{Fact}}{}
\@ifundefined{lem}{\newtheorem{lem}[thm]{Lemma}}{}
\@ifundefined{proofclaim}{\newtheorem{proofclaim}{Proofclaim}}{}
\makeatother
\makeatletter
\@ifundefined{proofclaim}{\newenvironment{proofclaim}[1][Proof of claim]{\begin{proof}[#1]\renewcommand*{\qedsymbol}{\(\blacksquare\)}}{\end{proof}}}{}
\makeatother
\title[Equitable tree colouring of graphs]{Equitable tree colouring of graphs}
\author{Yuping Gao, Allan Lo, Songling Shan}
\date{Source: arXiv:2604.13606v1 (CC BY 4.0)}
\begin{document}
\maketitle

\noindent\emph{Source and licence.} Equitable tree colouring of graphs. Version arXiv:2604.13606v1 (CC BY 4.0), distributed under the Creative Commons Attribution 4.0 International licence, as recorded in the arXiv licence metadata for this exact version and shown on the abstract page https://arxiv.org/abs/2604.13606v1. Full rights notice: licenses/arxiv-2604.13606-CC-BY-4.0.txt.\par\smallskip

\noindent\textit{Editorial scope.} Counted unit: the Claim whose source label is \texttt{clm:d=1:B} in the article (source0.tex lines 788--794, source label \texttt{clm:d=1:B}), delivered together with the article's own complete original proof of it (source0.tex lines 796--858, one \texttt{proof} environment of 63 lines). No mathematics is written, repaired or re-derived: statement and proof are the article's own text line for line. Required context retained uncounted: lemma:d-degenerate (lines 172--180); d-degenerate-a (lines 173--179); lemma:common-properties (lines 340--362); d (lines 344--351); f (lines 354--361). Every definition, display and label of those lines is retained unchanged. Omitted from this package and not redistributed: 1--171, 180--339, 352--353, 362--787, 795--795, 859--920. Nothing inside a retained excerpt is omitted or altered: every retained body is delivered exactly as the article's lines are written, so no omission receipt of any kind accompanies this package. Citations: the retained excerpts carry no citation command at all, so no bibliography entry of the article is transcribed and no reference is redistributed. Reference adaptation: none was needed and none was made; every reference inside the delivered unit resolves inside it, so no unresolved reference or citation is produced. Printed slips kept literal: no printed word, symbol, hypothesis or number of the counted statement or of its proof was altered. Layout and class: the article's own class is \texttt{article} with packages including geometry, babel, amssymb, amsmath, amsthm, bm, amsfonts, latexsym, lineno, enumitem, mathtools, hyperref; the wrapper is the lane's pinned \texttt{amsart} editorial wrapper at 10pt on A4, which restates the theorem-like environments the retained text needs and adds the article's own macro definitions that the retained text needs (none), with extra wrapper packages (bm, bbm, dsfont, xspace, cancel, mathtools, enumitem). The delivered numbering is the unit's own. Assets: no retained span contains a figure environment, a graphics-inclusion command, a PDF-page inclusion, a table or a TikZ picture; no image file is copied into this package, the package declares no asset dependency and no image file is redistributed with it. Licence: version arXiv:2604.13606v1 of this article is distributed under the Creative Commons Attribution 4.0 International licence, as recorded in the arXiv licence metadata for this exact version and shown on the abstract page https://arxiv.org/abs/2604.13606v1. A full rights notice is delivered at licenses/arxiv-2604.13606-CC-BY-4.0.txt. Characters: every retained excerpt is pure ASCII, so the delivered engine, XeLaTeX with its default Latin Modern text font, reports no missing character.
\begin{fact}\label{lemma:d-degenerate}~
	\begin{enumerate}[label={\rm(\alph*)}, start=1]
		\item  If $G'$ is obtained from $G$ by joining a new vertex to  at most $d$
		vertices in~$V^{*}$, then $G'$ is $d$-degenerate.\label{d-degenerate-a}
		\item 	 	If $G'$  is obtained
		from $G$ by joining a new vertex to  at most $d-1$
		vertices  of $G$, then $(V(G'))^{*}=V^{*}$.\label{d-degenerate-b}
\end{enumerate}
\end{fact}

	\begin{lem}\label{lemma:common-properties}
	Let $\Delta, d,k \in \mathbb{N}$ with $k \ge (\Delta+1)/(d+1)$, and $G$ be an $n$-vertex graph with $\Delta(G)\leq\Delta$.
		Suppose that $G$ has no equitable $d$-degenerate $k$-colouring, but $G$ has a minimal near-equitable $d$-degenerate $k$-colouring~$\varphi$.
	Then the following statements hold:
\begin{enumerate}[label={\rm (\alph*)},start=1]
\item  for all $V\in \mathcal{A}$,  $|V|-|C_-| \le 1$ and $V^*\neq \emptyset$; \label{a}
\item  for all $V\in \mathcal{A}$, $|B| \ge b |V|+1$;
 \label{b}		
\item for all $V\in \mathcal{T}$, each $v \in V$ is non-movable to any colour class in $\mathcal{A}\setminus ( \mathcal{T}\cup \{U_-\})$, in particular, $d_G(v,A \setminus (T \cup U_-)) \ge (d+1)(a-1-t)$; \label{c}
\item for all distinct $V,V'\in \mathcal{T}$, $\mathcal{D}-V$ contains a directed $(V',C_-)$-path; \label{d}
\item there exists $W\in \mathcal{T}$ such that for all $w\in W^{*}$ and $V\in \mathcal{T}$, we have $d_G(w,V) \ge d$.\label{e}	
\end{enumerate}			
Furthermore,  suppose that for all $x \in B$, $G[B]-x$ has an equitable $d$-degenerate $b$-colouring.
Then the following statements hold:	
\begin{enumerate}	[label={\rm (\alph*)},start=6]		
\item
if $v \in V^*$, $V  \in \mathcal{T}$ and $x \in N_G(v) \cap B$ such that $d_G(x,V^*) =d+1$, then $v$ is non-movable to any other colour class in~$\mathcal{A}$, in particular, $d_G(v,A) \ge (d+1)a-1$;
	  \label{f}
\item   if $ v \in V \in \mathcal{T}$ and there exist distinct $x,y \in B\cap N_{G}(v)$ such that
$d_G(x,V)  = d_G(y,V) =d+1$, then $x \sim y$ in~$G$; \label{g}		
\item   if $\tau \le d+1$ and $ k \ge \frac{1+\tau}{\tau} \frac{\Delta+1}{d+2}$, then $t < \tau b$. \label{h}
\end{enumerate}
\end{lem}

\begin{claim} \label{clm:d=1:B}
The following inequalities hold.
\begin{enumerate}[label={\rm(\alph*)},start=1]
	\item $2|B_1| \le e_G(B,W_1)$. \label{itm:d=1:B1}
	\item $2|B_1| + |B_2| \le e_G(B,W_1) +\frac13 e_G(B, W_2)$.\label{itm:d=1:B2}
\end{enumerate}
\end{claim}

\begin{proofclaim}
Observe that the following two inequalities imply the claim.
\begin{align}
		2|B_1| + 2 |B_{2,1}| &\le e_G(B,W_1), \label{eqn:d=1:B_1}\\
	2|B_1| + 2 |B_{2,1}| + 3|B_{2,2}| & \le e_G(B,W_1)+e_G(B,W_2). \label{eqn:d=1:B_2}
\end{align}
By Lemma~\ref{lemma:common-properties}\ref{f},  $N_G(B_1, W^*) \subseteq W_1$.
Together with the definition of~$B_{2,1}$, we have
\begin{align*}
	2|B_1| + 2 |B_{2,1}| \le \sum_{y \in B_1 \cup B_{2,1} } d_{G}(y,W_1) \le e_G(B,W_1)
\end{align*}
implying~\eqref{eqn:d=1:B_1}.

Suppose that $N_G(B_{2,2}, W^*) \subseteq W_1 \cup W_2$.
Thus,
\begin{align*}
	2|B_1| + 2 |B_{2,1}| + 3|B_{2,2}| & \le \sum_{y \in B_1 \cup B_{2,1} } d_{G}(y,W_1) +  \sum_{y \in B_{2,2} } d_{G}(y,W_1 \cup W_2)
	 \le e_G(B,W_1)+e_G(B,W_2)
\end{align*}
implying~\eqref{eqn:d=1:B_2}.

Therefore, we may assume that there exists $x \in B_{2,2}$ with $N_G(x,W_{>2}) \ne \emptyset$.
By our induction hypothesis, $G[B]-x$ has an equitable $1$-degenerate $b$-colouring~$\psi_B$.
Recall that $G$ does not have an equitable $1$-degenerate $k$-colouring.
Thus to obtain a contradiction (and to prove the claim), it suffices to show that $G[A \cup \{x\}]$ contains a $1$-degenerate $a$-colouring~$\psi_A$ in which the size of colour class~$C_-$ increases by one and the size of each remaining colour class in~$\mathcal{A}$ remains the same.


Since $x \notin B_{2,1}$, there exist distinct vertices $w_1 \in N_G(x,W_2 \cup W_{>2})$ and $w_2 \in N_G(x, W_{>2})$.
Since $w_1 \notin W_1$, there exists $V_1 \in \mathcal{T} \cup \{U_-\}$ such that $d(w_1, V^*_1) \le 1$, so $w_1$ is movable to~$V_1$.
By Lemma~\ref{lemma:common-properties}\ref{d}, there exists a directed $(V_1, C_-)$-path~$\mathcal{P}$ in~$\mathcal{D}-W$.
Note that $U_- \in V(\mathcal{P})$.
Similarly, since $w_2 \in W_{>2}$, there exist $U_1,U_2 \in \mathcal{R}'$ such that $w_2$ is movable to both $U_1$ and~$U_2$.



\noindent\textbf{Case 1: $U_1, U_2 \in V(\mathcal{P})$.}
Let $U_1$ be closer to~$V_1$ along~$\mathcal{P}$, so $V(V_1\mathcal{P}U_1) \cap V(U_2 \mathcal{P}C_-) = \emptyset$ and $U_1 \in \mathcal{R}$.
Since  $U_1 \in \mathcal{R}$, we have $e_G(U_1, W^*) < 2 (|U_1| -  \Delta^2)$.
There exist $\Delta^2+1$ vertices $u \in U_1$ such that $d_G(u, W^*) \le 1 $.
On the other hand,  there are at most $\Delta^2$ vertices that is joined to~$x$ with a path of length at most~$2$ and~$x \notin U_1$.
Thus, we can pick $u_1 \in U_1$ be such that $d(u_1, W^*) \le 1 $ and there is no path of length at most~$2$ between $u_1$ and~$x$, that is,
\begin{align}
	(\{x\} \cup N_G(x) ) \cap N_G(u_1) = \emptyset. \label{eqn:N(x)N(u1)}
\end{align}

We move $w_1$ to~$V_1$, a vertex of~$V_1$ to~$U_1$ along~$V_1\mathcal{P}U_1$, $w_2$ to~$U_2$, a vertex of~$U_2$ to~$C_-$ along~$U_2 \mathcal{P}C_-$ and both $x$ and $u_1$ to~$W$. Denote by~$\psi_A$ the resulting $a$-colouring of~$G[A \cup \{x\}]$.
The size of each colour class of~$\psi_A$ is as desired.

It remains to show that $\psi_A$ is indeed a $1$-degenerate colouring of~$G[A \cup \{x\}]$.
Each colour class~$V'$ in~$\psi_A$ not corresponding to~$W$ is either the same as before or is obtained by adding a movable vertex and deleting at most one vertex, so $G[V']$ is $1$-degenerate.
The colour class in~$\psi_A$ corresponding to~$W$ is $W_0 = W \cup \{x,u_1\} \setminus \{w_1,w_2\}$.
By Fact~\ref{lemma:d-degenerate}\ref{d-degenerate-a}, $G[W_0 \setminus \{u_1\}]$ is $1$-degenerate.
Observe that $(W_0 \setminus \{u_1\})^* \setminus W^*$ is a subset of isolated vertices in $G[W]$ and so $(W_0 \setminus \{u_1\})^* \setminus W^* \subseteq N_G(x)$.
By~\eqref{eqn:N(x)N(u1)}, we have $d_G(u_1, (W_0 \setminus \{u_1\})^*) \le d_G(u_1, W^*) \le 1$.
By Fact~\ref{lemma:d-degenerate}\ref{d-degenerate-a}, $G[W_0 ]$ is $1$-degenerate as required.


\noindent\textbf{Case 2: $U_i \notin V(\mathcal{P})$ for some $i \in [2]$.}
Since $U_- \in V(\mathcal{P})$, $U_i \in \mathcal{R}$ and so $e_G(U_i, W^*) < 2 (|U_i| -  \Delta^2)$.
By a similar argument as above there exists $u \in U_i$ such that $d_G(u, W^*) \le 1 $ and $(\{x\} \cup	N_G(x)) \cap N_G(u) = \emptyset$.
We move $w_1$ to~$V_1$, a vertex~$V_1$ to~$C_-$ along~$\mathcal{P}$, $w_2$ to~$U_i$ and both $x$ and $u$ to~$W$.
Again, by a similar argument, this results in a desired $1$-degenerate $a$-colouring of~$G[A \cup \{x\}]$.
\end{proofclaim}

\end{document}
